\documentclass[11pt]{article}
\usepackage[a4paper,margin=1in]{geometry}
\usepackage{amsmath,amssymb,mathtools,bm}
\usepackage{enumitem,booktabs,array}
\usepackage{tcolorbox}
\usepackage{hyperref}
\usepackage[protrusion=true,expansion=false]{microtype}
\hypersetup{colorlinks=true,linkcolor=blue,urlcolor=blue,citecolor=blue}
\setlist[itemize]{topsep=2pt,itemsep=2pt}
\setlist[enumerate]{topsep=2pt,itemsep=2pt}
\newtcolorbox{keybox}{colback=blue!3!white,colframe=blue!40!black,boxrule=0.4pt,arc=1mm}
\newtcolorbox{alertbox}{colback=red!3!white,colframe=red!40!black,boxrule=0.4pt,arc=1mm}
\newtcolorbox{answerbox}{colback=green!3!white,colframe=green!40!black,boxrule=0.4pt,arc=1mm}
\newtcolorbox{drillbox}{colback=yellow!5!white,colframe=orange!60!black,boxrule=0.4pt,arc=1mm}
\newcommand{\EV}{\mathbb{E}}
\newcommand{\Var}{\mathrm{Var}}
\newcommand{\Cov}{\mathrm{Cov}}
\newcommand{\blank}[1]{\underline{\hspace{#1}}}

\title{E11-04: DeAngelo, Goncalves \& Stulz (2022, RF) \\
\large ``Leverage and Cash Dynamics'' --- Active Recall Workbook}
\author{Empirical Corporate Finance II --- Exam Prep}
\date{\today}
\begin{document}\maketitle

\begin{keybox}
\textbf{Paper in one sentence.} DGS document the joint dynamics of leverage and cash holdings in US Compustat firms and argue for a \emph{net-debt} framework: firms manage gross debt and cash simultaneously, using cash as a flexibility buffer and debt for tax/discipline reasons; leverage studies that ignore cash misread capital-structure behavior --- many ``deleveraging'' episodes are actually just cash accumulation.

\textbf{Placement.} Modern re-framing of capital structure; consolidates insights from cash-holdings (Bates-Kahle-Stulz 2009), trade-off theory, and debt-structure literature. Important for joint capital-structure / liquidity-management research.
\end{keybox}

\section*{Round 1: Complete Walk-Through}

\subsection*{1.1 Motivation}
Corporate leverage measured as debt/assets has risen modestly over decades, but gross debt and cash have both grown. Whether firms have ``deleveraged'' depends on whether one looks at gross debt, net debt (debt minus cash), or total liabilities.

\subsection*{1.2 Framework: gross vs.\ net debt}
\begin{itemize}
\item Gross debt ratio: $D/A$.
\item Cash ratio: $C/A$.
\item Net debt ratio: $(D-C)/A$.
\item Capital-structure decisions are meaningful only if debt and cash are modeled jointly.
\end{itemize}

\subsection*{1.3 Data}
\begin{itemize}
\item Compustat US non-financial firms, 1980--2019.
\item Aggregate and firm-level gross debt, cash, and net debt ratios.
\item Decompose leverage changes into debt-issuance, debt-retirement, and cash-flow channels.
\end{itemize}

\subsection*{1.4 Headline facts}
\begin{itemize}
\item Average gross debt / assets rises modestly; average cash / assets rises sharply (especially post-1990s).
\item Net debt / assets is relatively stable over time.
\item Cross-sectionally, firms with higher precautionary-savings motives have lower net debt; firms with stronger tax benefits have higher net debt.
\item Firms that appear to ``deleverage'' are often just building cash reserves without reducing gross debt.
\end{itemize}

\subsection*{1.5 Mechanism}
\begin{itemize}
\item Cash is a flexibility / precautionary-savings instrument.
\item Gross debt gives tax shields and discipline; cash gives financial-flexibility option.
\item Optimal joint management balances these across macro and firm-level conditions.
\end{itemize}

\subsection*{1.6 Implications}
\begin{itemize}
\item Capital-structure research should use net debt or jointly model debt and cash.
\item Many ``deleveraging'' narratives misread data.
\item Supports dynamic trade-off theories with cash-management (e.g., Bolton-Chen-Wang 2011).
\end{itemize}

\subsection*{1.7 Caveats}
\begin{alertbox}
\begin{itemize}
\item Cash may include marketable securities with debt-like properties.
\item Net-debt formalism may over-smooth meaningful gross-debt changes.
\item Aggregate trends may mask cross-sectional heterogeneity.
\end{itemize}
\end{alertbox}

\section*{Round 2: Level 1 Fill-in}
\begin{enumerate}
\item Framework: gross debt, cash, \blank{2cm} debt.
\item Trend: gross \blank{2cm} rises modestly; cash rises sharply.
\item Net debt ratio: roughly \blank{2cm} over time.
\item Implication: ``deleveraging'' often is cash \blank{2cm}.
\item Related theory: \blank{2cm}-Chen-Wang 2011.
\end{enumerate}

\section*{Round 3: Level 2 Prompts}
\begin{enumerate}
\item Define gross, cash, and net debt.
\item Summarize secular trends.
\item Explain why ignoring cash biases leverage interpretation.
\item Relate DGS to Bates-Kahle-Stulz 2009.
\item Implications for capital-structure theory and applied work.
\end{enumerate}

\section*{Round 4: Minimal Scaffolding}
From a blank page: (i) framework, (ii) trends, (iii) mechanism, (iv) implications, (v) cautions.

\section*{Round 5: Blank Page}
Essay: joint dynamics of leverage and cash --- capital structure revisited.

\vspace{6pt}
\begin{drillbox}
\textbf{Speed Drill.} (1) Three variables. (2) Trend in cash. (3) Trend in net debt. (4) Reinterpretation. (5) Related cash-holdings paper.
\end{drillbox}

\section*{Quick Reference \& Answer Key}
\begin{answerbox}
\begin{itemize}
\item \textbf{Framework.} Gross debt + cash = net debt.
\item \textbf{Trend.} Cash up, net debt stable.
\item \textbf{Insight.} ``Deleveraging'' often is cash accumulation.
\item \textbf{Lesson.} Use net debt in capital-structure work.
\end{itemize}
\end{answerbox}

\begin{keybox}
\textbf{30-second exam pitch.} DeAngelo, Goncalves \& Stulz (2022) argue that empirical capital-structure research should treat cash and debt jointly. Using Compustat non-financial firms 1980--2019, they document that gross debt / assets has risen modestly while cash / assets has risen sharply (especially post-1990s), leaving net debt / assets relatively stable. Narratives of ``corporate deleveraging'' often describe cash accumulation, not reductions in gross debt. Cross-sectionally, firms facing stronger precautionary-savings motives run lower net debt, and those with higher tax benefits run higher net debt. The paper consolidates insights from Bates-Kahle-Stulz (2009) on cash trends with dynamic trade-off theory (Bolton-Chen-Wang 2011), arguing that capital-structure tests should use net debt or jointly-model gross debt and cash. It reshapes applied capital-structure work and clarifies secular trends in corporate balance sheets.
\end{keybox}

\end{document}
